% Electron-electron two-stream instability as a test of PIC numerics -- final report
% Builds on Overleaf as-is (pdfLaTeX + BibTeX), or locally with
% `latexmk twostream_electron_electron.tex`. Companion: twostream_electron_ion.tex.
% Missing figures render as labelled placeholder boxes; see figures/README.md.
% Audience: an undergraduate taking a first course in plasma physics.
\documentclass[10pt,twocolumn,letterpaper]{article}

\usepackage[letterpaper,margin=0.75in,columnsep=0.25in]{geometry}
\usepackage[T1]{fontenc}
\usepackage[utf8]{inputenc}
\usepackage{lmodern}
\usepackage{microtype}
\usepackage{amsmath,amssymb}
\usepackage{bm}
\usepackage{graphicx}
\usepackage{xcolor}
\usepackage{booktabs}
\usepackage{siunitx}
\usepackage[font=small,labelfont=bf]{caption}
\usepackage[numbers,sort&compress]{natbib}
\usepackage[colorlinks=true,allcolors=blue!50!black]{hyperref}

% ---- Notation -------------------------------------------------------------
\newcommand{\wb}{\omega_b}            % beam plasma frequency
\newcommand{\wpe}{\omega_{p}}         % total plasma frequency, wpe^2 = 2 wb^2
\newcommand{\vz}{v_0}                 % beam drift speed
\newcommand{\ppc}{N_\mathrm{ppc}}     % particles per cell
\newcommand{\dd}{\mathrm{d}}

% ---- Figure placeholder: \figorbox{file}{width}{label shown while missing} --
\newcommand{\figorbox}[3]{%
  \IfFileExists{figures/#1}%
    {\includegraphics[width=#2]{figures/#1}}%
    {\fbox{\parbox[c][0.55#2][c]{\dimexpr#2-2\fboxsep-2\fboxrule\relax}{\centering\small\textit{Placeholder:}\\ #3\\ \texttt{figures/\detokenize{#1}}}}}}

% ---- Margin notes for drafting (set \drafttrue -> \draftfalse to hide) -----
\newif\ifdraft \drafttrue
\newcommand{\todo}[1]{\ifdraft{\color{red}\textbf{[TODO:} #1\textbf{]}}\fi}

\title{\vspace{-1.5em}The Electron Two-Stream Instability as a Test of\\
       Particle-in-Cell Numerics}
\author{Harry Wu}
\date{\today}

\begin{document}
\maketitle

\begin{abstract}
\noindent
When two electron beams pass through each other, a small density ripple can grow exponentially
into a large electric wave: the two-stream instability. For cold, equal beams, theory predicts
the growth rate exactly, with maximum $\gamma_{\max}=\wb/2$, where $\wb$ is one beam's plasma
frequency. We reproduce the instability with the particle-in-cell (PIC) code \textsc{Entity}
and compare the measured growth rate with theory. Because the answer is known, the comparison
shows whether the simulation, with its approximations, reproduces the analytic result.
\todo{one sentence on the headline result once runs are in.}
\end{abstract}

% ===========================================================================
\section{Introduction}
\label{sec:intro}

Picture two streams of electrons moving through each other in opposite directions over a
neutralizing ion background. Suppose a small density ripple appears. Its electric field speeds
up some electrons and slows others, so they bunch up further where the density was already high.
More bunching gives a stronger field, which gives more bunching, and the ripple grows
exponentially until the wave is strong enough to trap electrons in its potential
wells~\citep{BohmGross1949,Chen2016}. A disturbance that grows on its own like this is an
\emph{instability}.

For cold beams (all electrons in a beam share one velocity) of equal density, the growth rate
can be calculated exactly, which makes this a standard test for computer
simulations~\citep{BirdsallLangdon1991}. A particle-in-cell (PIC) code cannot follow every
electron in a real plasma. It follows a smaller number of ``macroparticles'' and computes the
field on a grid. Both shortcuts are approximations: too few particles make the simulation noisy,
and the grid can artificially heat the plasma~\citep{Langdon1970,Okuda1972}. Here we show that,
with suitable settings, a PIC simulation reproduces the exact growth rate.

% ===========================================================================
\section{Linear theory}
\label{sec:theory}

Two cold electron beams, each of density $n_b$, move along $x$ at $\pm\vz$ through a fixed,
uniform ion background. Each beam has plasma frequency $\wb^2 = n_b e^2/(\varepsilon_0 m_e)$,
the frequency at which its electrons would oscillate on their own; for all electrons together
$\wpe^2=2\wb^2$. Treating each beam $j$ as a fluid, the continuity equation, Newton's second law,
and Gauss's law are
\begin{gather}
  \partial_t n_j + \partial_x (n_j v_j) = 0, \qquad
  m_e(\partial_t v_j + v_j\partial_x v_j) = -eE, \nonumber\\
  \varepsilon_0\,\partial_x E = -e\,(n_1 + n_2 - 2n_b).
\end{gather}
We write each quantity as its uniform value plus a small ripple, e.g.\
$n_j = n_b + \tilde n_j\, e^{i(kx-\omega t)}$, and keep only first-order terms. A beam moving
at $v_{0j}=\pm\vz$ sees the wave at the Doppler-shifted frequency $\omega - k v_{0j}$, and
\begin{equation}
  \tilde v_j = \frac{-ie\tilde E}{m_e(\omega - k v_{0j})}, \qquad
  \tilde n_j = \frac{k n_b\,\tilde v_j}{\omega - k v_{0j}}.
\end{equation}
Putting both $\tilde n_j$ into Gauss's law and dividing by $ik\tilde E$ gives the
\emph{dispersion relation}, which says which $\omega$ go with each $k$:
\begin{equation}
  1 = \frac{\wb^2}{(\omega - k\vz)^2} + \frac{\wb^2}{(\omega + k\vz)^2}.
  \label{eq:disp}
\end{equation}
This is a quadratic in $\omega^2$, with roots
$\omega^2 = k^2\vz^2 + \wb^2 \pm \wb\sqrt{\wb^2 + 4k^2\vz^2}$. For $k\vz<\sqrt2\,\wb$ the minus
root is negative, so $\omega = i\gamma$ and the ripple grows as $e^{\gamma t}$ with
\begin{equation}
  \gamma^2(k) = \wb\sqrt{\wb^2 + 4k^2\vz^2} - k^2\vz^2 - \wb^2 .
  \label{eq:gamma}
\end{equation}
Shorter waves (larger $k$) do not grow. Setting $\dd\gamma^2/\dd(k^2)=0$ gives the
fastest-growing wave,
\begin{equation}
  \gamma_{\max} = \tfrac{1}{2}\,\wb \quad\text{at}\quad k_{\max}\vz = \tfrac{\sqrt{3}}{2}\,\wb .
  \label{eq:gmax}
\end{equation}
In one growth time $1/\gamma$ the amplitude grows by a factor $e$ (one ``$e$-folding'').
Equation~\eqref{eq:gamma} is our reference curve. It holds only while the ripples are small
(hence \emph{linear} theory); once the wave traps electrons, growth stops and the electrons
form swirling vortices in phase space~\citep{BirdsallLangdon1991}.

% ===========================================================================
\section{The particle-in-cell method}
\label{sec:pic}

\paragraph{Macroparticles and noise.}
Each macroparticle stands for many real electrons, with their combined charge and mass. Its
$q/m$ is the same as an electron's, so it follows the same path. Fields are stored on grid
points spaced $\Delta x$ apart, with $\ppc$ particles per cell on average. With so few particles
the density is noisy, $\delta n/n\sim\ppc^{-1/2}$. The noise conveniently supplies the initial
ripple, but because growth starts from a larger amplitude, it saturates sooner and leaves fewer
$e$-foldings to measure~\citep{BirdsallLangdon1991}.

\paragraph{The PIC cycle.}
Each time step $\Delta t$ has four stages~\citep{BirdsallLangdon1991,HockneyEastwood1988}:
(1) \emph{gather} the field at each particle from nearby grid points; (2) \emph{push} the
particles; (3) \emph{deposit} their current onto the grid; (4) \emph{solve} for the new fields.
Steps 1 and 3 use a particle \emph{shape}. The simplest, nearest grid point, gives all of a
particle's charge to the closest grid point. Cloud-in-cell splits it between the two nearest
points according to distance, and higher orders spread it smoothly over three or four points.
Smoother shapes average out more noise but cost more per particle. This run uses cloud-in-cell.

\paragraph{Push.}
\textsc{Entity} uses the Boris method~\citep{Boris1970}. With no magnetic field and motion only
along $x$, it reduces to the \emph{leapfrog} scheme
\begin{equation}
  \begin{aligned}
    v^{n+1/2} &= v^{n-1/2} + \tfrac{q}{m}E^{n}\Delta t,\\
    x^{n+1} &= x^{n} + v^{n+1/2}\Delta t,
  \end{aligned}
\end{equation}
with velocities offset by half a step from positions for extra accuracy. To follow the plasma
oscillation the step must be short: leapfrog is stable for $\wpe\Delta t<2$ and accurate for
$\wpe\Delta t\lesssim0.2$.

\paragraph{Current and fields.}
The current is deposited with the zigzag scheme~\citep{Umeda2003}, which conserves charge
exactly, so Gauss's law
stays satisfied automatically, and it is smoothed (\emph{filtered}) over neighboring cells to
reduce noise. Fields are advanced on a staggered Yee grid~\citep{Yee1966}, which requires
$c\Delta t<\Delta x$. In our 1D problem Amp\`ere's law reduces to
$\partial_t E_x=-J_x/\varepsilon_0$, so the problem is effectively electrostatic.

\paragraph{Numerical heating.}
PIC codes do not conserve energy exactly. On a grid, waves whose wavenumbers differ by a
multiple of $2\pi/\Delta x$ look identical at every grid point (\emph{aliasing}), so
short-wavelength noise can pose as longer waves. When the Debye length is much smaller than a
cell, $\lambda_D\ll\Delta x$, this heats the plasma until $\lambda_D\sim\Delta x$: the
finite-grid instability~\citep{Langdon1970,Okuda1972}. Heating would matter here, because a
velocity spread $v_{th}$ smears out the bunching and lowers $\gamma$ once $k v_{th}\sim\gamma$, so
we choose a grid fine enough to avoid it (Sec.~\ref{sec:setup}).

% ===========================================================================
\section{Simulation setup}
\label{sec:setup}

We use \textsc{Entity}~\citep{Entity} (v1.3.3, \texttt{streaming} problem generator, input
\texttt{twostream.toml}) in 1D with periodic boundaries. Code units have $c=1$ and a reference
skin depth $d_0=c/\omega_{p0}=10$, so $\omega_{p0}=0.1$; the physics depends only on ratios such
as $\gamma/\wb$. There are four species: the two electron beams, and two positron species that
are never moved and act as the fixed ion background. The beams have four-velocity
$u_x=\gamma_L v_x/c=\pm0.1$, where $\gamma_L\approx1$ is the Lorentz factor, so they are
non-relativistic.

The input \texttt{densities = [0.5, 0.5]} sets the density of each electron--positron pair, so
$n_b=0.25\,n_0$ and $\wb=\omega_{p0}\sqrt{0.25}=0.05$. Equation~\eqref{eq:gmax} then predicts
$\gamma_{\max}=0.025$ at $k_{\max}\approx0.435$ ($\lambda_{\max}\approx14.4$). A periodic box
of length $L=100$ allows only $k=2\pi m/L$; the $m=7$ mode has $k\vz/\wb=0.875$, close to the
peak at $0.866$. The run lasts $\wb t=50$, or 25 growth times $1/\gamma_{\max}$. Because it
starts from random particle noise (no seeded wave), the fastest wave grows only about 3
$e$-folds before it saturates; the rest of the run shows the nonlinear stage. Table~\ref{tab:params} lists all parameters. The Debye length of
all the electrons together is $\lambda_D=v_{th}/\wpe=0.14$, about 17 cells ($\Delta x/\lambda_D\approx0.06$),
far from the regime $\lambda_D\ll\Delta x$ where finite-grid heating occurs.
\todo{update if the production runs change any of these values.}

\begin{table}[t]
  \centering
  \caption{Simulation parameters from \texttt{twostream.toml} in code units ($c=1$), with
  dimensionless ratios on the right.}
  \label{tab:params}
  \small
  \begin{tabular}{@{}lll@{}}
    \toprule
    Parameter & Value & Normalized \\
    \midrule
    Beam four-velocity $u_x$            & $\pm0.1$            & $\vz\approx0.0995\,c$ \\
    Beam density $n_b$                  & $0.25\,n_0$         & $\wb=0.05$ \\
    Temperature $T/m_ec^2$              & $10^{-4}$           & $v_{th}/\vz\approx0.1$ \\
    Box length $L$                      & 100                 & $L\wb/\vz\approx50$ \\
    Cells                               & 12288               & $\Delta x/\lambda_D\approx0.06$ \\
    Time step (CFL $=0.5$)              & $\Delta t\approx0.0041$ & $\wb\Delta t\approx2\times10^{-4}$ \\
    Particles per cell $\ppc$           & 16                  & 4 per species \\
    Current filter passes               & 4                   & \\
    Run time                            & 1000                & $\wb t=50$ \\
    Particle shape                      & cloud-in-cell       & first order \\
    Pusher, deposit                     & Boris, zigzag       & \\
    \bottomrule
  \end{tabular}
\end{table}

\paragraph{Measuring $\gamma(k)$.}
We Fourier transform $E_x(x,t)$ in $x$ to get the amplitude $|\tilde E_k(t)|$ of each box wave.
During growth $|\tilde E_k|\propto e^{\gamma t}$, so $\ln|\tilde E_k|$ versus $t$ is a straight
line of slope $\gamma(k)$. We fit it between the noise floor and saturation, and take the spread
over several fit windows as the uncertainty (Fig.~\ref{fig:growth}).

\begin{figure}[t]
  \centering
  \figorbox{ee_growth_mode.pdf}{\columnwidth}{$\ln|\tilde E_k(t)|$ for $k\approx k_{\max}$, with fit}
  \caption{Amplitude of the fastest-growing wave versus time. The dashed line is the
  straight-line fit during exponential growth; its slope is $\gamma$. \todo{caption details.}}
  \label{fig:growth}
\end{figure}

% ===========================================================================
\section{Results}
\label{sec:results}

\paragraph{Growth rate versus wavenumber.}
Figure~\ref{fig:gammak} compares the measured $\gamma(k)$ with Eq.~\eqref{eq:gamma}. For the
fastest wave ($m=7$, $k\vz/\wb=0.875$) we measure $\gamma=0.51\pm0.17\,\wb$
($0.36\pm0.12\,\wpe$), against the theory value $\wb/2$ ($0.354\,\wpe$). The uncertainty is
large because moving the fit window within $t=40$--$160$ changes the result by
$0.06$--$0.19\,\wpe$. Averaged over the well-resolved waves $m=4$--$9$, the ratio of measured to
theoretical growth rate is $0.99\pm0.16$ (per-wave ratios 0.85, 0.96, 0.93, 1.01, 1.23, 0.98), so
on average the simulation is consistent with the cold-beam theory. The uncertainty propagates
each wave's fit-window spread. We quote it instead of the smaller scatter-based standard error
($0.05$) because all six waves grow from the same noise realization and so are not independent. Individual waves scatter by about $\pm30\%$ ($\pm25$--$50\%$). The scatter
comes from the random-noise start: with only about 3 $e$-folds of clean growth, each fit is
short. The longest waves ($m=1$--$3$) barely rise above the noise and are uninformative, and the
wave nearest the cutoff $k\vz=\sqrt2\,\wb$ ($m=11$) sits high because the strongly growing
modes drive it. A more precise test would seed the wave or average several runs with different
random seeds. Both need code changes: the \texttt{streaming} problem generator loads only uniform
drifting Maxwellians, and \textsc{Entity}'s random seed is a compile-time constant
(\texttt{constant::RandomSeed}), not an input parameter.

\begin{figure}[t]
  \centering
  \figorbox{ee_gamma_vs_k.pdf}{\columnwidth}{measured $\gamma(k)/\wb$ vs.\ $k\vz/\wb$ with theory}
  \caption{Measured growth rate of each box wave (points) compared with the cold-beam theory,
  Eq.~\eqref{eq:gamma} (line). Error bars show the spread of $\gamma$ over fit windows within
  $t=40$--$160$. The well-resolved waves $m=4$--$9$ agree with theory on average
  (ratio $0.99\pm0.16$) but scatter individually by about $\pm30\%$; $m=1$--$3$ are
  noise-dominated and $m=11$ is driven by the fastest modes.}
  \label{fig:gammak}
\end{figure}

% ===========================================================================
\section{Discussion and conclusions}
\label{sec:discussion}

\todo{Summarize: (i) PIC recovers $\gamma(k)$ on average (ratio $0.99\pm0.16$ for $m=4$--$9$),
with $\sim\pm30\%$ scatter per wave; (ii) limitations: random-noise start gives only $\sim3$
$e$-folds of growth, the small thermal spread lowers $\gamma$ slightly, and only box modes are
sampled; (iii) seeding or averaging over random seeds would sharpen the test, but both need code
changes (pgen initialization; hard-coded \texttt{constant::RandomSeed}).}

{\small\setlength{\bibsep}{1pt}
\bibliographystyle{unsrtnat}
\bibliography{refs}}

\end{document}
